% ============================================================================
% CHAPITRE 4 : ÉTUDE COMPLÈTE DE FONCTIONS
% ============================================================================
\chapter{Étude Complète de Fonctions}

La \textbf{Question 1} de l'examen d'admission d'Analyse à l'UCLouvain (qui représente typiquement $15$ points sur $50$, soit $30\%$ de la note totale !) est systématiquement consacrée à l'\textbf{étude complète d'une fonction} $f : \D \subset \R \to \R$ aboutissant à son tableau de variations et à l'esquisse précise de son graphe sur un repère quadrillé. Ce chapitre expose la théorie reliant les dérivées $f'$ et $f''$ à la géométrie de la courbe (croissance, extrema locaux, concavité, points d'inflexion), puis détaille pas à pas la méthodologie attendue par les correcteurs à travers cinq études complètes issues des annales 2021 à 2026.

% ----------------------------------------------------------------------------
\section{Croissance, extrema locaux, concavité et points d'inflexion}
% ----------------------------------------------------------------------------

\subsection{Signe de la dérivée première $f'$ et sens de variation}

En appliquant le théorème des accroissements finis $f(x_2) - f(x_1) = f'(c)(x_2 - x_1)$ entre deux points quelconques $x_1 < x_2$ d'un intervalle $I$, on voit immédiatement que le signe de $f(x_2) - f(x_1)$ est le même que celui de $f'(c)$.

\begin{theoreme}[Théorème fondamental des variations sur un intervalle $I$]
Soit $f$ une fonction continue sur un intervalle $I$ et dérivable à l'intérieur de $I$ :
\begin{enumerate}
    \item $f'(x) \ge 0$ sur $I \iff f$ est \textbf{croissante} sur $I$ ($\nearrow$).\\
    Si $f'(x) > 0$ sur $I$ (sauf éventuellement en un nombre fini de points isolés où $f'(x) = 0$), alors $f$ est \textbf{strictement croissante} sur $I$.
    \item $f'(x) \le 0$ sur $I \iff f$ est \textbf{décroissante} sur $I$ ($\searrow$).\\
    Si $f'(x) < 0$ sur $I$ (sauf éventuellement en des points isolés), alors $f$ est \textbf{strictement décroissante} sur $I$.
    \item $f'(x) = 0$ pour tout $x \in I \iff f$ est \textbf{constante} sur l'intervalle $I$.
\end{enumerate}
\end{theoreme}

\subsection{Extrema locaux (minima et maxima) : conditions du 1\ier{} et du 2\ieme{} ordre}

\begin{definition}[Maximum local, minimum local et extremum global]
Soit $f : \D \subset \R \to \R$ et soit $x^* \in \D$ :
\begin{itemize}
    \item On dit que $f$ admet un \textbf{maximum local} (ou relatif) en $x^*$ s'il existe un intervalle ouvert $]x^* - \delta, x^* + \delta[$ ($\delta > 0$) tel que $\forall x \in \D \cap \,]x^* - \delta, x^* + \delta[,\; f(x) \le f(x^*)$. Si cette inégalité est vraie sur tout le domaine $\D$, c'est un \textbf{maximum global} (ou absolu).
    \item On dit que $f$ admet un \textbf{minimum local} (ou relatif) en $x^*$ s'il existe $\delta > 0$ tel que $\forall x \in \D \cap \,]x^* - \delta, x^* + \delta[,\; f(x) \ge f(x^*)$ (et \textbf{minimum global} si c'est vrai sur tout $\D$).
    \item Un \textbf{point critique} (ou point stationnaire) de $f$ est un point $x^* \in \D$ où la dérivée s'annule : $\boxed{f'(x^*) = 0}$.
\end{itemize}
\end{definition}

\begin{theoreme}[Caractérisation des extrema locaux d'une fonction dérivable]
Soit $f$ une fonction dérivable sur un intervalle ouvert $I \subset \D$ et soit $x^* \in I$ un point intérieur à $I$ :
\begin{enumerate}
    \item \textbf{Condition nécessaire du 1\ier{} ordre (Théorème de Fermat) :}
    \[
        \boxed{\text{Si } f \text{ a un extremum local en un point intérieur } x^* \in I, \text{ alors } f'(x^*) = 0}.
    \]
    \item \textbf{Condition suffisante du 1\ier{} ordre (changement de signe de $f'$) :} Soit $x^* \in I$ tel que $f'(x^*) = 0$ :
    \begin{itemize}
        \item Si $f'(x)$ passe du signe \textbf{$+$ à gauche au signe $-$ à droite} de $x^*$ ($\nearrow \;\text{max}\; \searrow$), alors $f$ admet un \textbf{maximum local} en $x^*$.
        \item Si $f'(x)$ passe du signe \textbf{$-$ à gauche au signe $+$ à droite} de $x^*$ ($\searrow \;\text{min}\; \nearrow$), alors $f$ admet un \textbf{minimum local} en $x^*$.
        \item Si $f'(x)$ \textbf{garde le même signe} à gauche et à droite de $x^*$ (tout en s'annulant en $x^*$, comme $x^3$ en $0$ ou la fonction de \textbf{Septembre 2026 Q1} en $x = 1$), alors $x^*$ \textbf{n'est PAS un extremum local} : c'est un \textbf{point d'inflexion à tangente horizontale} !
    \end{itemize}
    \item \textbf{Condition suffisante du 2\ieme{} ordre (signe de la dérivée seconde $f''(x^*)$) :} Si $f$ est deux fois dérivable en un point critique $x^* \in I$ ($f'(x^*) = 0$) :
    \[
        \boxed{\begin{cases}
            f'(x^*) = 0 \;\text{ et }\; f''(x^*) < 0 &\implies f \text{ admet un \textbf{maximum local strict} en } x^* \;\; (\cap) \\
            f'(x^*) = 0 \;\text{ et }\; f''(x^*) > 0 &\implies f \text{ admet un \textbf{minimum local strict} en } x^* \;\; (\cup)
        \end{cases}}
    \]
    \emph{Attention :} Si $f'(x^*) = 0$ et $f''(x^*) = 0$, le test de la dérivée seconde seul ne permet pas de conclure : il faut regarder si $f'(x)$ change de signe de part et d'autre de $x^*$ (par exemple, pour $x^4$ ou $f(x) = x^2\sin(x^2)+\cos(x^2)$ en $x=0$ en Juillet 2023 Série 1 Q1, on a $f'(0) = f''(0) = 0$ et c'est un minimum très plat ; tandis que pour $x^3$ en $0$, $f'(0) = f''(0) = 0$ et c'est un point d'inflexion !).
\end{enumerate}
\end{theoreme}

\subsection{Signe de la dérivée seconde $f''$, concavité et points d'inflexion}

\begin{definition}[Concavité tournée vers le haut / vers le bas et point d'inflexion]
Soit $f$ une fonction deux fois dérivable sur un intervalle $I$ :
\begin{enumerate}
    \item \textbf{Concavité positive (tournée vers le haut $\cup$, fonction convexe) :}
    \[
        \boxed{f''(x) > 0 \text{ sur } I \implies f' \text{ croissante} \implies \text{concavité tournée vers le haut } (\cup)}.
    \]
    Géométriquement, la courbe est située \textbf{au-dessus de toutes ses tangentes} sur $I$.
    \item \textbf{Concavité négative (tournée vers le bas $\cap$, fonction concave) :}
    \[
        \boxed{f''(x) < 0 \text{ sur } I \implies f' \text{ décroissante} \implies \text{concavité tournée vers le bas } (\cap)}.
    \]
    Géométriquement, la courbe est située \textbf{en dessous de toutes ses tangentes} sur $I$.
    \item \textbf{Point d'inflexion (PI) :} Un point $x_0 \in \D$ (où $f$ est continue) est un \textbf{point d'inflexion} du graphe de $f$ si la fonction \textbf{change de sens de concavité} en passant par $x_0$ (c'est-à-dire si la dérivée seconde $f''(x)$ \textbf{s'annule en $x_0$ en changeant de signe} de part et d'autre de $x_0$). En un point d'inflexion, la courbe \textbf{traverse sa droite tangente}.
\end{enumerate}
\end{definition}

\begin{center}
\begin{tikzpicture}[scale=0.88, >=Stealth]
    % Concavité vers le haut (minimum)
    \begin{scope}[shift={(0,0)}]
        \draw[very thick, uclblue, domain=-1.5:1.5, samples=60] plot (\x, {0.65*(\x)^2 + 0.4});
        \draw[thick, defgreen] (-1.2,0.4) -- (1.2,0.4);
        \fill[defgreen] (0,0.4) circle (2pt) node[below] {\footnotesize Min ($f'=0, f''>0$)};
        \node[above, uclblue, font=\bfseries\footnotesize] at (0,2.1) {Concavité positive ($f'' > 0, \cup$)};
        \node[below, darkgray, font=\footnotesize] at (0,-0.25) {Au-dessus des tangentes};
    \end{scope}

    % Concavité vers le bas (maximum)
    \begin{scope}[shift={(5.0,0)}]
        \draw[very thick, warnred, domain=-1.5:1.5, samples=60] plot (\x, {1.85 - 0.65*(\x)^2});
        \draw[thick, defgreen] (-1.2,1.85) -- (1.2,1.85);
        \fill[defgreen] (0,1.85) circle (2pt) node[above] {\footnotesize Max ($f'=0, f''<0$)};
        \node[above, warnred, font=\bfseries\footnotesize] at (0,2.4) {Concavité négative ($f'' < 0, \cap$)};
        \node[below, darkgray, font=\footnotesize] at (0,-0.25) {En dessous des tangentes};
    \end{scope}

    % Point d'inflexion
    \begin{scope}[shift={(10.0,0.9)}]
        \draw[very thick, exorange, domain=-1.45:1.45, samples=60] plot (\x, {0.42*(\x)^3 + 0.15*\x});
        \draw[thick, dashed, darkgray] (-1.3,-0.195) -- (1.3,0.195);
        \fill[exorange] (0,0) circle (2.2pt) node[below right] {\footnotesize PI ($f''=0$)};
        \node[above, exorange, font=\bfseries\footnotesize] at (0,1.5) {Point d'inflexion ($\cap \to \cup$)};
        \node[below, darkgray, font=\footnotesize] at (0,-1.15) {Traverse sa tangente};
    \end{scope}
\end{tikzpicture}
\end{center}

\begin{exemple}[Extrema locaux en fonction de paramètres : UCLouvain -- Juillet 2021, Question 3a]
\emph{Énoncé :} Soient $a, b, c$ des constantes non nulles. On considère la fonction $f$ définie sur $\R_0$ par :
\[
    f(x) = \frac{ax^2 + bx + c}{x^2}.
\]
En quel(s) point(s) $x$ la fonction $f$ admet-elle un maximum et/ou un minimum local ? S'agit-il d'un maximum ou d'un minimum (la réponse peut dépendre des valeurs de $a, b, c$) ?

\emph{Résolution :}
\begin{itemize}
    \item \emph{Astuce de réécriture avant de dériver :} Au lieu d'appliquer la lourde formule de dérivation d'un quotient $\frac{u}{v}$, divisons chaque terme du numérateur par $x^2$ :
    \[
        f(x) = a + \frac{b}{x} + \frac{c}{x^2} = a + b\,x^{-1} + c\,x^{-2} \quad (\forall x \neq 0).
    \]
    \item \emph{Dérivée première et point critique :} En dérivant terme à terme :
    \[
        f'(x) = -\frac{b}{x^2} - \frac{2c}{x^3} = -\frac{bx + 2c}{x^3}.
    \]
    Comme $b \neq 0$ et $c \neq 0$, la dérivée s'annule en un unique point critique $x^* \neq 0$ :
    \[
        f'(x^*) = 0 \iff bx^* + 2c = 0 \iff \boxed{x^* = -\frac{2c}{b}}.
    \]
    \item \emph{Nature de l'extremum par la dérivée seconde $f''(x^*)$ :} Dérivons une seconde fois $f'(x) = -b\,x^{-2} - 2c\,x^{-3}$ :
    \[
        f''(x) = \frac{2b}{x^3} + \frac{6c}{x^4} = \frac{2bx + 6c}{x^4}.
    \]
    Évaluons $f''$ en $x^* = -\frac{2c}{b}$ (où $2bx^* + 6c = 2b\left(-\frac{2c}{b}\right) + 6c = -4c + 6c = 2c$) :
    \[
        f''(x^*) = \frac{2c}{(x^*)^4} = \frac{2c}{\left(-\frac{2c}{b}\right)^4} = \frac{2c\,b^4}{16c^4} = \frac{b^4}{8c^3}.
    \]
    Puisque $b^4 > 0$ pour tout $b \neq 0$, le signe de $f''(x^*)$ ne dépend que du signe de $c$ (ni $a$ ni $b$ n'influencent la nature de l'extremum !) :
    \begin{itemize}
        \item \textbf{Si $c > 0$ :} $f''(x^*) > 0$, donc $f$ admet un \textbf{minimum local} en $x^* = -\dfrac{2c}{b}$.
        \item \textbf{Si $c < 0$ :} $f''(x^*) < 0$, donc $f$ admet un \textbf{maximum local} en $x^* = -\dfrac{2c}{b}$.
    \end{itemize}
    \emph{Remarque élégante du solutionnaire :} En posant $y = \frac{1}{x}$, on a $f(x) = a + by + cy^2$, qui est une simple parabole en $y$ dont le sommet est en $y^* = -\frac{b}{2c} \iff x^* = -\frac{2c}{b}$, tournée vers le haut (minimum) si $c > 0$ et vers le bas (maximum) si $c < 0$ !
\end{itemize}
\end{exemple}

% ----------------------------------------------------------------------------
\section{Méthodologie en 6 étapes pour réussir la Question 1 de l'examen}
% ----------------------------------------------------------------------------

À l'examen d'admission de l'UCLouvain, la Question 1 est guidée par des sous-questions a), b), c), d), e) (et parfois f)). Voici la feuille de route méthodologique qui répond exactement aux attentes du barème officiel :

\begin{astuce}[Les 6 étapes du succès à la Question 1 (Étude de fonction)]
\begin{enumerate}
    \item \textbf{Domaine $\D$, symétries, racines et comportement aux bords :}
    \begin{itemize}
        \item Écrire les conditions d'existence et exprimer clairement $\D$.
        \item Repérer si $f$ est paire, impaire ou périodique.
        \item Chercher les racines $f(x) = 0$ si l'équation est résoluble, ainsi que les limites aux bords ouverts de $\D$ (asymptotes verticales AV, horizontales AH ou obliques AO).
    \end{itemize}
    \item \textbf{Dérivée première $f'(x)$ et ses racines :}
    \begin{itemize}
        \item Calculer $f'(x)$ proprement et \textbf{la mettre sous forme factorisée / sur un même dénominateur} (sans factorisation, il est impossible d'étudier son signe !).
        \item Résoudre $f'(x) = 0$ sur $\D$ pour identifier les candidats extrema locaux.
    \end{itemize}
    \item \textbf{Dérivée seconde $f''(x)$ et ses racines (si demandée) :}
    \begin{itemize}
        \item Partir de la forme la plus simple de $f'(x)$, calculer $f''(x)$ et la factoriser pour trouver ses racines (candidats points d'inflexion) et son signe.
    \end{itemize}
    \item \textbf{Calcul des valeurs exactes ET numériques approchées aux points remarquables :}
    \begin{itemize}
        \item Depuis 2024, l'énoncé fournit une \textbf{Aide numérique} (par ex. $\sqrt{2} \simeq 1,414$, $\sqrt{5} \simeq 2,236$, $e^{-1} \simeq 0,368$, $e^{-2} \simeq 0,135$, $\ln(2) \simeq 0,693$, $\pi^2 \simeq 9,87$).
        \item Calculez au brouillon, puis reportez dans le tableau la \textbf{valeur mathématique exacte ET la valeur décimale approchée} de $f(x)$ (et de $f'(x)$ si demandé) en chaque point remarquable : c'est indispensable pour placer les points au bon endroit sur le quadrillage !
    \end{itemize}
    \item \textbf{Tableau de variations dans l'ordre logique officiel :}
    \begin{itemize}
        \item Les énoncés récents (2025--2026) demandent explicitement de présenter les lignes du tableau dans l'\textbf{ordre logique} :
        \[
            \mathbf{1.\; x} \;\to\; \mathbf{2.\; f'(x)} \;\text{(signe)} \;\to\; \mathbf{3.\; f''(x)} \;\text{(signe)} \;\to\; \mathbf{4.\; f(x)} \;\text{(valeurs, } \nearrow\searrow, \cup\cap\text{)}.
        \]
        \item Indiquez dans la dernière ligne les mentions \textbf{max}, \textbf{min}, \textbf{inflexion (PI)}, \textbf{racine} et \textbf{AV}.
    \end{itemize}
    \item \textbf{Esquisse graphique sur le quadrillage avec axes chiffrés :}
    \begin{itemize}
        \item \textbf{Graduez et chiffrez vos deux axes $x$ et $y$} en choisissant une échelle adaptée aux valeurs de votre tableau.
        \item Tracez d'abord les \textbf{asymptotes en pointillés} (AV, AH, AO), puis placez tous les \textbf{points remarquables} (racines, extrema avec un petit segment horizontal pour la tangente horizontale, points d'inflexion, bords avec tangente verticale éventuelle).
        \item Reliez les points par une courbe lisse respectant scrupuleusement le sens de concavité ($\cup$ ou $\cap$) et annotez clairement sur le dessin : \emph{racine}, \emph{max}, \emph{min}, \emph{point d'inflexion}, \emph{AV}.
    \end{itemize}
\end{enumerate}
\end{astuce}

% ----------------------------------------------------------------------------
\section{Cinq études complètes de fonctions issues des annales UCLouvain}
% ----------------------------------------------------------------------------

Étudions maintenant de A à Z cinq sujets représentatifs couvrant toute la diversité des fonctions posées à l'examen d'admission entre 2022 et 2026.

\begin{exemple}[Étude 1 : Fonction trigonométrique périodique -- UCLouvain, Septembre 2022, Question 1]
\emph{Énoncé :} On considère la fonction $f : \D \subset \R \to \R$ définie par :
\[
    f(x) = \frac{\cos^2(x)}{1 + \cos(x)}.
\]
\begin{enumerate}[label=\alph*)]
    \item Quel est le domaine de définition $\D$ de $f$ ?
    \item Calculez la fonction dérivée $f'$.
    \item Quels sont les minima et maxima locaux éventuels de $f$ ? \emph{(Note : il ne vous est pas demandé de calculer la dérivée seconde).}
    \item Sur base de vos résultats précédents, esquissez le graphique de la fonction $f$ pour $x \in [-\pi, \pi]$. Veillez à aussi y indiquer la position des racines, s'il y en a.
\end{enumerate}

\emph{Résolution complète :}
\begin{enumerate}[label=\alph*)]
    \item \textbf{Domaine, parité, périodicité et racines :}
    \begin{itemize}
        \item \textbf{CE :} $1 + \cos(x) \neq 0 \iff \cos(x) \neq -1 \iff x \neq \pi + 2k\pi$ ($k \in \Z$). D'où $\D = \R \setminus \{\pi + 2k\pi \mid k \in \Z\}$.
        \item La fonction $f$ est \textbf{paire} ($f(-x) = f(x)$) et \textbf{périodique de période $2\pi$}. Sur $[-\pi, \pi]$, le domaine est $]-\pi, \pi[$.
        \item \textbf{Racines :} $f(x) = 0 \iff \cos^2(x) = 0$ (avec $\cos(x) \neq -1$) $\iff \cos(x) = 0 \iff x = \frac{\pi}{2} + k\pi$. Sur $]-\pi, \pi[$, il y a deux racines : $x = -\frac{\pi}{2}$ et $x = \frac{\pi}{2}$.
        \item \textbf{Asymptotes verticales :} Pour $x \to \pm\pi$ (à l'intérieur de $]-\pi, \pi[$), $\cos(x) \to -1^+$ donc $1 + \cos(x) \to 0^+$ et $\cos^2(x) \to (-1)^2 = 1 > 0$. Ainsi $\lim_{x \to \pi^-} f(x) = \lim_{x \to -\pi^+} f(x) = \frac{1}{0^+} = +\infty$, d'où les asymptotes verticales $\text{AV} \equiv x = \pi + 2k\pi$.
    \end{itemize}
    \item \textbf{Dérivée première $f'(x)$ :} Par la règle du quotient avec $(\cos^2 x)' = -2\cos(x)\sin(x)$ :
    \begin{align*}
        f'(x) &= \frac{-2\cos(x)\sin(x)\bigl(1+\cos(x)\bigr) - \cos^2(x)\bigl(-\sin(x)\bigr)}{\bigl(1+\cos(x)\bigr)^2} \\
        &= \boxed{-\frac{\cos(x)\sin(x)\bigl(2 + \cos(x)\bigr)}{\bigl(1+\cos(x)\bigr)^2}} = -\frac{\sin(2x)\bigl(2+\cos(x)\bigr)}{2\bigl(1+\cos(x)\bigr)^2}.
    \end{align*}
    \item \textbf{Signe de $f'(x)$ et extrema locaux :} Sur tout le domaine $\D$, on a $\bigl(1+\cos(x)\bigr)^2 > 0$ et $2 + \cos(x) \ge 2 - 1 = 1 > 0$. Le signe de $f'(x)$ est donc exactement le signe de $-\cos(x)\sin(x) = -\frac{1}{2}\sin(2x)$.
    \[
        f'(x) = 0 \iff \sin(x) = 0 \text{ ou } \cos(x) = 0 \iff x = 2k\pi \text{ ou } x = \frac{\pi}{2} + k\pi \;\; (k \in \Z).
    \]
    Dressons le tableau de signes et de variations sur $[-\pi, \pi]$ :
    \begin{center}
    {\small
    \renewcommand{\arraystretch}{1.35}
    \setlength{\tabcolsep}{4pt}
    \begin{tabular}{|c||c|c|c|c|c|c|c|c|c|}
    \hline
    $x$ & $-\pi$ & & $-\dfrac{\pi}{2}$ & & $0$ & & $\dfrac{\pi}{2}$ & & $\pi$ \\ \hline\hline
    $\sin(2x)$ & $0$ & $+$ & $0$ & $-$ & $0$ & $+$ & $0$ & $-$ & $0$ \\ \hline
    $1 + \cos(x)$ & $0$ & $+$ & $1$ & $+$ & $2$ & $+$ & $1$ & $+$ & $0$ \\ \hline
    $f'(x)$ & AV & $-$ & $0$ & $+$ & $0$ & $-$ & $0$ & $+$ & AV \\ \hline
    $f(x)$ & $\text{AV}\;(+\infty)$ & $\searrow$ & $\mathbf{0}$ (min) & $\nearrow$ & $\mathbf{\dfrac{1}{2}}$ (max) & $\searrow$ & $\mathbf{0}$ (min) & $\nearrow$ & $(+\infty)\;\text{AV}$ \\ \hline
    \end{tabular}
    }
    \end{center}
    La fonction admet donc :
    \begin{itemize}
        \item des \textbf{minima locaux} (et globaux, qui sont aussi les racines de $f$, de valeur $0$) en $\boxed{x = \dfrac{\pi}{2} + k\pi}$ ($k \in \Z$) ;
        \item des \textbf{maxima locaux} (de valeur $f(0) = \dfrac{1}{2}$) en $\boxed{x = 2k\pi}$ ($k \in \Z$).
    \end{itemize}
    \item \textbf{Esquisse graphique sur $[-\pi, \pi]$ :}
    \begin{center}
    \begin{tikzpicture}[scale=1.05, >=Stealth]
        \begin{axis}[
            width=10.5cm, height=6.5cm,
            axis lines=middle,
            xmin=-3.6, xmax=3.6,
            ymin=-0.3, ymax=3.2,
            xtick={-3.1416, -1.5708, 0, 1.5708, 3.1416},
            xticklabels={$-\pi$, $-\frac{\pi}{2}$, $0$, $\frac{\pi}{2}$, $\pi$},
            ytick={0.5, 1, 2, 3},
            yticklabels={$\frac{1}{2}$, $1$, $2$, $3$},
            xlabel={$x$},
            ylabel={$f(x)$},
            grid=both,
            grid style={gray!20},
            samples=120
        ]
            \addplot[very thick, uclblue, domain=-2.72:2.72] {(cos(deg(x)))^2 / (1 + cos(deg(x)))};
            \draw[dashed, warnred, thick] (axis cs:-3.1416, -0.2) -- (axis cs:-3.1416, 3.1) node[pos=0.85, right] {\small AV};
            \draw[dashed, warnred, thick] (axis cs:3.1416, -0.2) -- (axis cs:3.1416, 3.1) node[pos=0.85, left] {\small AV};
            \fill[defgreen] (axis cs:-1.5708, 0) circle (2.2pt) node[below] {\footnotesize min \& racine};
            \fill[defgreen] (axis cs:1.5708, 0) circle (2.2pt) node[below] {\footnotesize min \& racine};
            \fill[exorange] (axis cs:0, 0.5) circle (2.2pt) node[above right] {\footnotesize max $(0, 1/2)$};
        \end{axis}
    \end{tikzpicture}
    \end{center}
\end{enumerate}
\end{exemple}

\begin{exemple}[Étude 2 : Fonction algébrique avec radical -- UCLouvain, Juillet 2024, Question 1]
\emph{Énoncé :} On considère la fonction $f : \D \subset \R \to \R$ définie par :
\[
    f(x) = \sqrt{1 + x} - x.
\]
\begin{enumerate}[label=\alph*)]
    \item Quel est le domaine de définition $\D$ de $f$ ? Obtenez aussi la(les) racine(s) de $f$.
    \item Obtenez la fonction dérivée première $f'$.
    \item Obtenez la fonction dérivée seconde $f''$.
    \item Réalisez un tableau avec les points $x = -1, x = 0, x = 1, x = 2$, la(les) racine(s) de $f$ et le(s) point(s) où $f'$ s'annule. Reportez-y clairement toutes les valeurs de $f$ et $f'$, ainsi que le signe de $f''$.
    \item Esquissez proprement le graphe de $f$ dans l'intervalle $[-1, 2]$. \textbf{Aide :} $\sqrt{2} \simeq 1,414$, $\sqrt{3} \simeq 1,732$, $\sqrt{5} \simeq 2,236$.
\end{enumerate}

\emph{Résolution complète :}
\begin{enumerate}[label=\alph*)]
    \item \textbf{Domaine et racine :} Comme démontré à l'Exemple du Chapitre 1, $1 + x \ge 0 \iff x \ge -1$, donc $\boxed{\D = [-1, +\infty[}$. L'équation $\sqrt{1+x} = x$ (qui impose $x \ge 0$ et $x^2 - x - 1 = 0$) admet \textbf{une unique racine} dans $\D$ :
    \[
        \boxed{x_1 = \frac{1 + \sqrt{5}}{2} \simeq \frac{1 + 2,236}{2} = 1,618}.
    \]
    \item \textbf{Dérivée première $f'(x)$ et point critique :} Pour tout $x > -1$ :
    \[
        \boxed{f'(x) = \frac{1}{2\sqrt{1+x}} - 1}.
    \]
    (En $x \to -1^+$, $\lim_{x \to -1^+} f'(x) = +\infty$ : demi-tangente verticale au bord gauche $x = -1$). Cherchons où $f'$ s'annule :
    \[
        f'(x) = 0 \iff \sqrt{1+x} = \frac{1}{2} \iff 1 + x = \frac{1}{4} \iff \boxed{x = -\frac{3}{4} = -0,75}.
    \]
    En ce point $x = -\frac{3}{4}$, la fonction vaut $f\left(-\frac{3}{4}\right) = \sqrt{\frac{1}{4}} - \left(-\frac{3}{4}\right) = \frac{1}{2} + \frac{3}{4} = \dfrac{5}{4} = 1,25$.
    \item \textbf{Dérivée seconde $f''(x)$ :} En dérivant $f'(x) = \frac{1}{2}(1+x)^{-1/2} - 1$ pour $x > -1$ :
    \[
        \boxed{f''(x) = -\frac{1}{4}(1+x)^{-3/2} = -\frac{1}{4(1+x)\sqrt{1+x}} < 0 \quad (\forall x > -1)}.
    \]
    La dérivée seconde est \textbf{strictement négative partout} sur $]-1, +\infty[$ : la concavité est toujours tournée vers le bas ($\cap$), il n'y a \textbf{aucun point d'inflexion}, et le point critique $x = -\frac{3}{4}$ est un \textbf{maximum global}.
    \item \textbf{Tableau de variations et de valeurs sur $[-1, 2]$ :}
    \begin{center}
    {\scriptsize
    \renewcommand{\arraystretch}{1.45}
    \setlength{\tabcolsep}{1.4pt}
    \begin{tabular}{|c||c|c|c|c|c|c|c|c|c|c|c|}
    \hline
    $x$ & $-1$ & & $-\dfrac{3}{4}$ & & $0$ & & $1$ & & $x_1 = \dfrac{1+\sqrt{5}}{2}$ & & $2$ \\ \hline\hline
    $f'(x)$ & $+\infty$ & $+$ & $0$ & $-$ & $-\dfrac{1}{2}$ & $-$ & $\dfrac{1}{2\sqrt{2}}-1 \simeq -0,65$ & $-$ & $\dfrac{1}{2x_1}-1 < 0$ & $-$ & $\dfrac{1}{2\sqrt{3}}-1 \simeq -0,71$ \\ \hline
    $f''(x)$ & $-\infty$ & $-$ & $-$ & $-$ & $-$ & $-$ & $-$ & $-$ & $-$ & $-$ & $-$ \\ \hline
    $f(x)$ & $1$ & $\nearrow$ & $\mathbf{\dfrac{5}{4} = 1,25}$ (max) & $\searrow$ & $1$ & $\searrow$ & $\sqrt{2}-1 \simeq 0,41$ & $\searrow$ & $\mathbf{0}$ (racine) & $\searrow$ & $\sqrt{3}-2 \simeq -0,27$ \\ \hline
    \end{tabular}
    }
    \end{center}
    \item \textbf{Esquisse graphique sur $[-1, 2]$ :}
    \begin{center}
    \begin{tikzpicture}[scale=1.05, >=Stealth]
        \begin{axis}[
            width=10.5cm, height=6.5cm,
            axis lines=middle,
            xmin=-1.25, xmax=2.25,
            ymin=-0.45, ymax=1.45,
            xtick={-1, -0.75, 0, 1, 1.618, 2},
            xticklabels={$-1$, $-\frac{3}{4}$, $0$, $1$, $x_1 \simeq 1{,}62$, $2$},
            ytick={-0.268, 0.414, 1, 1.25},
            yticklabels={$-0{,}27$, $0{,}41$, $1$, $\frac{5}{4}$},
            xlabel={$x$},
            ylabel={$f(x)$},
            grid=both,
            grid style={gray!20},
            samples=150
        ]
            \addplot[very thick, uclblue, domain=-1:2.05] {sqrt(1+x) - x};
            \draw[thick, warnred, <->] (axis cs:-0.95, 1.25) -- (axis cs:-0.55, 1.25);
            \draw[thick, exorange, ->] (axis cs:-1, 1) -- (axis cs:-1, 1.32) node[above right] {\footnotesize pente $\infty$};
            \fill[exorange] (axis cs:-1, 1) circle (2pt);
            \fill[warnred] (axis cs:-0.75, 1.25) circle (2.2pt) node[above] {\footnotesize max $(-3/4, 5/4)$};
            \fill[uclblue] (axis cs:0, 1) circle (2pt);
            \fill[uclblue] (axis cs:1, 0.414) circle (2pt);
            \fill[defgreen] (axis cs:1.618, 0) circle (2.2pt) node[above right] {\footnotesize racine $x_1$};
            \fill[uclblue] (axis cs:2, -0.268) circle (2pt);
        \end{axis}
    \end{tikzpicture}
    \end{center}
\end{enumerate}
\end{exemple}

\begin{exemple}[Étude 3 : Fonction logarithmique et trigonométrique -- UCLouvain, Juillet 2025, Question 1]
\emph{Énoncé :} On considère la fonction $f : \D \subset \R \to \R$ définie par :
\[
    f(x) = x + \ln\bigl(\sin(x) + \cos(x)\bigr).
\]
\begin{enumerate}[label=\alph*)]
    \item Obtenez le domaine de définition $\D = \bigcup_{k\in\Z} \,]a_k, b_k[$.
    \item Obtenez la fonction dérivée première $f'$.
    \item Obtenez la fonction dérivée seconde $f''$.
    \item Réalisez le tableau de variation de $f$ pour la partie de $\D$ comprise dans $[-\pi, \pi]$, dans l'ordre logique $f'$, puis $f''$, puis $f$.
    \item Esquissez proprement le graphe de $f$ sur cette même partie de $\D$.
\end{enumerate}

\emph{Résolution complète :}
\begin{enumerate}[label=\alph*)]
    \item \textbf{Domaine :} Vu à l'Exemple du Chapitre 1 : $\sin(x) + \cos(x) = \sqrt{2}\cos\left(x - \frac{\pi}{4}\right) > 0 \iff x \in \,]-\frac{\pi}{4} + 2k\pi,\; \frac{3\pi}{4} + 2k\pi[$. Dans $[-\pi, \pi]$, la partie du domaine est l'intervalle ouvert $\left]-\dfrac{\pi}{4}, \dfrac{3\pi}{4}\right[$.
    \item \textbf{Dérivée première $f'(x)$ :} En dérivant terme à terme et en mettant au même dénominateur :
    \[
        f'(x) = 1 + \frac{\cos(x) - \sin(x)}{\sin(x) + \cos(x)} = \frac{\sin(x) + \cos(x) + \cos(x) - \sin(x)}{\sin(x) + \cos(x)} = \boxed{\frac{2\cos(x)}{\sin(x) + \cos(x)}}.
    \]
    Sur $\left]-\frac{\pi}{4}, \frac{3\pi}{4}\right[$, le dénominateur $\sin(x) + \cos(x)$ est strictement positif. Donc $f'(x)$ a le même signe que $\cos(x)$, et s'annule uniquement en $\cos(x) = 0 \iff \boxed{x = \dfrac{\pi}{2}}$ (avec $f'(x) > 0$ sur $\left]-\frac{\pi}{4}, \frac{\pi}{2}\right[$ et $f'(x) < 0$ sur $\left]\frac{\pi}{2}, \frac{3\pi}{4}\right[$).
    \item \textbf{Dérivée seconde $f''(x)$ :} Dérivons le quotient $f'(x) = \dfrac{2\cos(x)}{\sin(x)+\cos(x)}$ :
    \begin{align*}
        f''(x) &= 2\,\frac{-\sin(x)\bigl(\sin(x)+\cos(x)\bigr) - \cos(x)\bigl(\cos(x)-\sin(x)\bigr)}{\bigl(\sin(x)+\cos(x)\bigr)^2} \\
        &= \frac{2\bigl(-\sin^2(x) - \cos^2(x)\bigr)}{\bigl(\sin(x)+\cos(x)\bigr)^2} = \boxed{-\frac{2}{\bigl(\sin(x)+\cos(x)\bigr)^2} < 0}.
    \end{align*}
    (On peut aussi l'écrire $-\dfrac{2}{1 + \sin(2x)}$ car $(\sin x + \cos x)^2 = \sin^2 x + \cos^2 x + 2\sin x\cos x = 1 + \sin(2x)$).
    La dérivée seconde est \textbf{strictement négative partout} : la concavité est négative ($\cap$) sur tout l'intervalle, sans point d'inflexion.
    \item \textbf{Comportement aux bornes, racines et tableau de variations sur $\left]-\frac{\pi}{4}, \frac{3\pi}{4}\right[$ :}
    \begin{itemize}
        \item En $x \to -\frac{\pi}{4}^+$ et en $x \to \frac{3\pi}{4}^-$, $\sin(x) + \cos(x) \to 0^+ \implies \ln(\sin x + \cos x) \to -\infty$, d'où deux \textbf{asymptotes verticales} $\text{AV}_1 \equiv x = -\dfrac{\pi}{4}$ et $\text{AV}_2 \equiv x = \dfrac{3\pi}{4}$.
        \item Au maximum $x = \frac{\pi}{2}$, on a $f\left(\frac{\pi}{2}\right) = \frac{\pi}{2} + \ln(1 + 0) = \dfrac{\pi}{2} \simeq 1,57$.
        \item En $x = 0$, $f(0) = 0 + \ln(0 + 1) = 0$ : $x = 0$ est une première \textbf{racine} évidente (avec pente $f'(0) = \frac{2}{1} = 2$). Comme $f$ redescend continûment de $f\left(\frac{\pi}{2}\right) = \frac{\pi}{2} > 0$ vers $-\infty$ en $\frac{3\pi}{4}^-$, le TVI garantit une \textbf{seconde racine} dans $\left]\frac{\pi}{2}, \frac{3\pi}{4}\right[$.
    \end{itemize}
    \begin{center}
    \renewcommand{\arraystretch}{1.4}
    \begin{tabular}{|c||c|c|c|c|c|c|c|}
    \hline
    $x$ & $-\dfrac{\pi}{4}$ & & $0$ & & $\dfrac{\pi}{2}$ & & $\dfrac{3\pi}{4}$ \\ \hline\hline
    $f'(x)$ & AV & $+$ & $2$ & $+$ & $0$ & $-$ & AV \\ \hline
    $f''(x)$ & AV & $-$ & $-2$ & $-$ & $-2$ & $-$ & AV \\ \hline
    $f(x)$ & $\text{AV}\;(-\infty)$ & $\nearrow\;\cap$ & $\mathbf{0}$ (racine) & $\nearrow\;\cap$ & $\mathbf{\dfrac{\pi}{2}}$ (max) & $\searrow\;\cap$ & $(-\infty)\;\text{AV}$ \\ \hline
    \end{tabular}
    \end{center}
    \item \textbf{Esquisse graphique sur $\left]-\frac{\pi}{4}, \frac{3\pi}{4}\right[$ :}
    \begin{center}
    \begin{tikzpicture}[scale=1.05, >=Stealth]
        \begin{axis}[
            width=10.0cm, height=6.5cm,
            axis lines=middle,
            xmin=-1.1, xmax=2.65,
            ymin=-2.4, ymax=1.95,
            xtick={-0.7854, 0, 0.7854, 1.5708, 2.3562},
            xticklabels={$-\frac{\pi}{4}$, $0$, $\frac{\pi}{4}$, $\frac{\pi}{2}$, $\frac{3\pi}{4}$},
            ytick={-2, -1, 0.7854, 1.5708},
            yticklabels={$-2$, $-1$, $\frac{\pi}{4}$, $\frac{\pi}{2}$},
            xlabel={$x$},
            ylabel={$f(x)$},
            grid=both,
            grid style={gray!20},
            samples=140
        ]
            \addplot[very thick, uclblue, domain=-0.75:2.335] {x + ln(sin(deg(x)) + cos(deg(x)))};
            \draw[dashed, warnred, thick] (axis cs:-0.7854, -2.35) -- (axis cs:-0.7854, 1.85) node[pos=0.15, right] {\small AV};
            \draw[dashed, warnred, thick] (axis cs:2.3562, -2.35) -- (axis cs:2.3562, 1.85) node[pos=0.15, left] {\small AV};
            \draw[thick, warnred, <->] (axis cs:1.25, 1.5708) -- (axis cs:1.89, 1.5708);
            \fill[warnred] (axis cs:1.5708, 1.5708) circle (2.2pt) node[above] {\footnotesize max $(\pi/2, \pi/2)$};
            \fill[defgreen] (axis cs:0, 0) circle (2.2pt) node[below right] {\footnotesize racine $(0,0)$};
            \fill[defgreen] (axis cs:2.253, 0) circle (2.2pt) node[above left] {\footnotesize 2\ieme{} racine};
        \end{axis}
    \end{tikzpicture}
    \end{center}
\end{enumerate}
\end{exemple}

\begin{exemple}[Étude 4 : Logarithmes, extremum et point d'inflexion -- UCLouvain, Septembre 2025, Question 1]
\emph{Énoncé :} On considère la fonction $f : \D \subset \R \to \R$ définie par :
\[
    f(x) = x\,\ln^2(x).
\]
\textbf{Aide :} $\lim_{y \to 0^+} y^p\ln(y) = 0$ pour $p > 0$, $e \simeq 2,72$, $\dfrac{1}{e} \simeq 0,368$ et $\ln(2) \simeq 0,693$.
\begin{enumerate}[label=\alph*)]
    \item Obtenez le domaine de définition $\D$ et la/les racine(s) de $f$.
    \item Obtenez la dérivée première $f'$ et ses racines.
    \item Obtenez la dérivée seconde $f''$ et ses racines.
    \item Réalisez le tableau de variation de $f$ (ordre logique $f'$, $f''$, $f$).
    \item Esquissez proprement le graphe de $f$ pour $x \le 2$ en indiquant les valeurs numériques aux points remarquables.
\end{enumerate}

\emph{Résolution complète :}
\begin{enumerate}[label=\alph*)]
    \item \textbf{Domaine, comportement en $0^+$ et racines :}
    \begin{itemize}
        \item \textbf{CE :} $x > 0 \implies \boxed{\D = \,]0, +\infty[}$.
        \item \textbf{Limite en $0^+$ :} $\lim_{x \to 0^+} x\ln^2(x) = \lim_{x \to 0^+} \bigl(x^{1/2}\ln(x)\bigr)^2 = 0^2 = 0$ (pas d'asymptote verticale en $0$, la courbe part du point limite $(0, 0)$).
        \item \textbf{Racine sur $\D = \,]0, +\infty[$ :} $f(x) = 0 \iff \ln(x) = 0 \iff \boxed{x = 1}$ (et $f(x) > 0$ pour tout $x \in \D \setminus \{1\}$).
    \end{itemize}
    \item \textbf{Dérivée première $f'(x)$ et ses racines :} Par la règle du produit et de la chaîne $\bigl(\ln^2(x)\bigr)' = 2\ln(x)\cdot\frac{1}{x}$ :
    \[
        f'(x) = 1\cdot\ln^2(x) + x \cdot 2\ln(x)\frac{1}{x} = \ln^2(x) + 2\ln(x) = \boxed{\ln(x)\bigl(\ln(x) + 2\bigr)}.
    \]
    (Notons qu'en $x \to 0^+$, $\lim_{x \to 0^+} f'(x) = (-\infty)(-\infty) = +\infty$ : demi-tangente verticale en $0^+$).
    \[
        f'(x) = 0 \iff \ln(x) = 0 \;\text{ ou }\; \ln(x) = -2 \iff \boxed{x = 1 \;\text{ ou }\; x = e^{-2} = \frac{1}{e^2} \simeq 0,135}.
    \]
    \item \textbf{Dérivée seconde $f''(x)$ et ses racines :} En dérivant $f'(x) = \ln^2(x) + 2\ln(x)$ :
    \[
        f''(x) = \frac{2\ln(x)}{x} + \frac{2}{x} = \boxed{\frac{2}{x}\bigl(\ln(x) + 1\bigr)}.
    \]
    Comme $x > 0$, $f''(x)$ s'annule en changeant de signe uniquement en :
    \[
        f''(x) = 0 \iff \ln(x) = -1 \iff \boxed{x = e^{-1} = \frac{1}{e} \simeq 0,368} \quad (\textbf{point d'inflexion}).
    \]
    \item \textbf{Valeurs remarquables et tableau de variations :}
    \begin{itemize}
        \item En $x = e^{-2} \simeq 0,135$ (maximum local) : $f(e^{-2}) = e^{-2}(-2)^2 = \dfrac{4}{e^2} \simeq 4 \times 0,135 = \mathbf{0,542}$.
        \item En $x = e^{-1} \simeq 0,368$ (point d'inflexion) : $f(e^{-1}) = e^{-1}(-1)^2 = \dfrac{1}{e} \simeq \mathbf{0,368}$.
        \item En $x = 1$ (minimum local et racine) : $f(1) = \mathbf{0}$.
        \item En $x = 2$ (bord droit du tracé) : $f(2) = 2\ln^2(2) \simeq 2 \times (0,693)^2 \simeq 2 \times 0,480 = \mathbf{0,960}$.
    \end{itemize}
    \begin{center}
    {\footnotesize
    \renewcommand{\arraystretch}{1.4}
    \setlength{\tabcolsep}{2.8pt}
    \begin{tabular}{|c||c|c|c|c|c|c|c|c|}
    \hline
    $x$ & $0$ & & $e^{-2} \simeq 0,135$ & & $e^{-1} \simeq 0,368$ & & $1$ & \\ \hline\hline
    $f'(x)$ & $+\infty$ & $+$ & $0$ & $-$ & $-1$ & $-$ & $0$ & $+$ \\ \hline
    $f''(x)$ & $-\infty$ & $-$ & $-2e^2 < 0$ & $-$ & $0$ & $+$ & $2 > 0$ & $+$ \\ \hline
    $f(x)$ & $\lim = 0$ & $\nearrow\;\cap$ & $\mathbf{\dfrac{4}{e^2} \simeq 0,542}$ (max) & $\searrow\;\cap$ & $\mathbf{\dfrac{1}{e} \simeq 0,368}$ (PI) & $\searrow\;\cup$ & $\mathbf{0}$ (min \& racine) & $\nearrow\;\cup$ \\ \hline
    \end{tabular}
    }
    \end{center}
    \item \textbf{Esquisse graphique sur $]0, 2]$ :}
    \begin{center}
    \begin{tikzpicture}[scale=1.05, >=Stealth]
        \begin{axis}[
            width=10.0cm, height=6.5cm,
            axis lines=middle,
            xmin=-0.1, xmax=2.15,
            ymin=-0.12, ymax=1.08,
            xtick={0.135, 0.368, 1, 2},
            xticklabels={$0{,}135$, $0{,}368$, $1$, $2$},
            ytick={0.368, 0.541, 0.960},
            yticklabels={$0{,}368$, $0{,}542$, $0{,}960$},
            xlabel={$x$},
            ylabel={$f(x)$},
            grid=both,
            grid style={gray!20},
            samples=200
        ]
            \addplot[very thick, uclblue, domain=0.001:2.03] {x * (ln(x))^2};
            \draw[uclblue, thick] (axis cs:0,0) circle (2pt);
            \fill[warnred] (axis cs:0.1353, 0.5413) circle (2.2pt) node[above right] {\footnotesize max $(e^{-2}, 4/e^2)$};
            \fill[exorange] (axis cs:0.3679, 0.3679) circle (2.2pt) node[above right] {\footnotesize inflexion $(1/e, 1/e)$};
            \fill[defgreen] (axis cs:1, 0) circle (2.2pt) node[above right] {\footnotesize min \& racine $(1,0)$};
            \fill[uclblue] (axis cs:2, 0.961) circle (2pt) node[left] {\footnotesize $(2;\; 0{,}96)$};
        \end{axis}
    \end{tikzpicture}
    \end{center}
\end{enumerate}
\end{exemple}

\begin{exemple}[Étude 5 : Le cas d'un point d'inflexion à tangente horizontale -- UCLouvain, Septembre 2026, Question 1]
\emph{Énoncé :} On considère la fonction $f : \,]0, +\infty[ \;\to \R$ définie par :
\[
    f(x) = \frac{x^2}{2}\left(\ln^2(x) - \ln(x) + \frac{1}{2}\right).
\]
\textbf{Aide :} $\lim_{y \to 0^+} y^p\ln(y) = 0$ pour $p > 0$ ; $e^{-2} \simeq 0,135$ ; $f(2) \simeq 0,575$.
\begin{enumerate}[label=\alph*)]
    \item Obtenez $\lim_{x \to 0} f(x)$ en détail.
    \item Obtenez la dérivée première $f'$, sa/ses racine(s) ainsi que $\lim_{x \to 0} f'(x)$.
    \item Obtenez la dérivée seconde $f''$, sa/ses racine(s) ainsi que $\lim_{x \to 0} f''(x)$.
    \item Réalisez le tableau de variation de $f$ (avec valeurs numériques approchées).
    \item Esquissez proprement le graphe de $f$ pour $x \le 2$.
\end{enumerate}

\emph{Résolution complète :}
\begin{enumerate}[label=\alph*)]
    \item \textbf{Limite en $0^+$ :} En développant $f(x)$ :
    \[
        \lim_{x \to 0^+} f(x) = \frac{1}{2}\lim_{x \to 0^+} \bigl(x\ln(x)\bigr)^2 - \frac{1}{2}\lim_{x \to 0^+} \bigl(x^2\ln(x)\bigr) + \frac{1}{4}\lim_{x \to 0^+} x^2 = 0 - 0 + 0 = \boxed{0}.
    \]
    (De plus, comme le trinôme $u^2 - u + \frac{1}{2} = \left(u - \frac{1}{2}\right)^2 + \frac{1}{4} \ge \frac{1}{4} > 0$ ne s'annule jamais, $f$ ne possède \textbf{aucune racine} sur $]0, +\infty[$ !).
    \item \textbf{Dérivée première $f'(x)$ (simplification remarquable !) :} Par la règle du produit :
    \begin{align*}
        f'(x) &= x\left(\ln^2(x) - \ln(x) + \frac{1}{2}\right) + \frac{x^2}{2}\left(2\ln(x)\cdot\frac{1}{x} - \frac{1}{x}\right) \\
        &= x\ln^2(x) - \cancel{x\ln(x)} + \cancel{\frac{x}{2}} + \cancel{x\ln(x)} - \cancel{\frac{x}{2}} = \boxed{x\,\ln^2(x)}.
    \end{align*}
    \emph{Remarque :} La fonction $f(x)$ de Septembre 2026 est tout simplement la primitive (trouvée en Septembre 2025 Q3 !) de la fonction $x\ln^2(x)$ étudiée en Septembre 2025 Q1 !
    \begin{itemize}
        \item \textbf{Racine de $f'$ :} $f'(x) = x\ln^2(x) = 0 \iff \ln(x) = 0 \iff \boxed{x = 1}$.
        \item \textbf{Attention au signe de $f'(x)$ :} Sur $]0, +\infty[$, on a $x > 0$ et $\ln^2(x) \ge 0$ (c'est un carré !). Donc $f'(x) \ge 0$ \textbf{partout} : $f'(x)$ s'annule en $x = 1$ \textbf{sans changer de signe} ! La fonction $f$ est donc \textbf{strictement croissante sur tout $]0, +\infty[$} et n'admet \textbf{aucun extremum local} !
        \item \textbf{Limite de $f'$ en $0^+$ :} $\lim_{x \to 0^+} f'(x) = \lim_{x \to 0^+} \bigl(x^{1/2}\ln(x)\bigr)^2 = 0^2 = \boxed{0}$ (demi-tangente horizontale en $(0,0)$ !).
    \end{itemize}
    \item \textbf{Dérivée seconde $f''(x)$ :} Comme $f'(x) = x\ln^2(x)$, on retrouve le calcul de l'Étude 4 :
    \[
        \boxed{f''(x) = \ln^2(x) + 2\ln(x) = \ln(x)\bigl(\ln(x) + 2\bigr)}.
    \]
    \begin{itemize}
        \item \textbf{Racines de $f''$ (deux points d'inflexion !) :} $f''(x) = 0 \iff \ln(x) = -2$ ou $\ln(x) = 0 \iff \boxed{x = e^{-2} \simeq 0,135}$ ou $\boxed{x = 1}$. Comme $f''(x)$ change de signe en chacun de ces deux points ($+$ avant $e^{-2}$, $-$ entre $e^{-2}$ et $1$, $+$ après $1$), la courbe possède \textbf{deux points d'inflexion}, dont celui en $x = 1$ est un \textbf{point d'inflexion à tangente horizontale} ($f'(1) = f''(1) = 0$).
        \item \textbf{Limite de $f''$ en $0^+$ :} $\lim_{x \to 0^+} \ln(x)\bigl(\ln(x)+2\bigr) = (-\infty)(-\infty) = \boxed{+\infty}$.
    \end{itemize}
    \item \textbf{Valeurs remarquables et tableau de variations :}
    \begin{itemize}
        \item En $x = e^{-2} \simeq 0,135$ : $f'(e^{-2}) = 4e^{-2} \simeq 0,540$ et $f(e^{-2}) = \frac{e^{-4}}{2}\left(4 - (-2) + \frac{1}{2}\right) = \frac{13}{4}e^{-4} \simeq 3,25 \times (0,135)^2 \simeq \mathbf{0,059}$.
        \item En $x = 1$ : $f'(1) = 0$, $f''(1) = 0$ et $f(1) = \frac{1}{2}\left(0 - 0 + \frac{1}{2}\right) = \dfrac{1}{4} = \mathbf{0,250}$.
    \end{itemize}
    \begin{center}
    {\small
    \renewcommand{\arraystretch}{1.4}
    \setlength{\tabcolsep}{4pt}
    \begin{tabular}{|c||c|c|c|c|c|c|}
    \hline
    $x$ & $\text{limite } 0$ & & $e^{-2} \simeq 0,135$ & & $1$ & \\ \hline\hline
    $f'(x)$ & $0$ & $+$ & $4e^{-2} \simeq 0,540$ & $+$ & $0$ & $+$ \\ \hline
    $f''(x)$ & $+\infty$ & $+$ & $0$ & $-$ & $0$ & $+$ \\ \hline
    $f(x)$ & $0$ & $\nearrow\;\cup$ & $\mathbf{\dfrac{13}{4}e^{-4} \simeq 0,059}$ (PI) & $\nearrow\;\cap$ & $\mathbf{\dfrac{1}{4} = 0,250}$ (PI tang. horiz.) & $\nearrow\;\cup$ \\ \hline
    \end{tabular}
    }
    \end{center}
    \item \textbf{Esquisse graphique sur $]0, 2]$ :}
    \begin{center}
    \begin{tikzpicture}[scale=1.05, >=Stealth]
        \begin{axis}[
            width=10.0cm, height=6.5cm,
            axis lines=middle,
            xmin=-0.1, xmax=2.15,
            ymin=-0.05, ymax=0.65,
            xtick={0.135, 1, 2},
            xticklabels={$0{,}135$, $1$, $2$},
            ytick={0.059, 0.250, 0.500, 0.575},
            yticklabels={$0{,}059$, $0{,}250$, $0{,}500$, $0{,}575$},
            xlabel={$x$},
            ylabel={$f(x)$},
            grid=both,
            grid style={gray!20},
            samples=200
        ]
            \addplot[very thick, uclblue, domain=0.001:2.03] {0.5 * x^2 * ((ln(x))^2 - ln(x) + 0.5)};
            \draw[thick, warnred, <->] (axis cs:0.65, 0.25) -- (axis cs:1.35, 0.25);
            \fill[exorange] (axis cs:0.1353, 0.0595) circle (2.2pt) node[below right] {\footnotesize PI $(e^{-2};\; 0{,}059)$};
            \fill[warnred] (axis cs:1, 0.25) circle (2.2pt) node[above left] {\footnotesize PI à tangente horizontale $(1;\; 0{,}25)$};
            \fill[uclblue] (axis cs:2, 0.575) circle (2pt) node[left] {\footnotesize $(2;\; 0{,}575)$};
        \end{axis}
    \end{tikzpicture}
    \end{center}
\end{enumerate}
\end{exemple}

% ----------------------------------------------------------------------------
\section{Exercices du Chapitre 4}
% ----------------------------------------------------------------------------

\emph{Les corrigés détaillés de tous ces exercices se trouvent au Chapitre 9, en fin d'ouvrage.}

\begin{enumerate}[label=\textbf{Exercice 4.\arabic* --}]
    \item \textbf{(UCLouvain -- Juillet 2021, Q1 -- Étude complète)} On considère la fonction $f : \D \subset \R \to \R$ définie par :
    \[
        f(x) = \ln\left(\frac{1}{\cos(x)}\right).
    \]
    \begin{enumerate}
        \item Déterminer le domaine de définition $\D$ de $f$.
        \item Préciser le domaine de dérivabilité et calculer $f'(x)$ ainsi que la dérivée seconde $f''(x)$.
        \item Déterminer les minima et maxima locaux éventuels de $f$ ainsi que ses asymptotes.
        \item Dresser le tableau de variations et de concavité sur $[-\pi, \pi[$, puis expliquer comment en déduire le graphe sur $[-4\pi, 4\pi]$ de $f(x)$ et de $\bar{f}(x) = \ln\left(\dfrac{1}{\cos^2(x)}\right)$.
    \end{enumerate}

    \item \textbf{(UCLouvain -- Juillet 2022, Q1 -- Étude complète)} Réaliser l'étude complète (domaine $\D$, dérivée première $f'$, dérivée seconde $f''$, extrema locaux, asymptotes, racines et tableau de variations sur $[-\pi, \pi]$) de la fonction :
    \[
        f(x) = \ln\left(\frac{1}{1 + \cos(x)}\right).
    \]

    \item \textbf{(UCLouvain -- Juillet 2023 -- Série 1, Q1 -- Étude complète)} On considère la fonction $f : \R \to \R$ définie par :
    \[
        f(x) = x^2\sin(x^2) + \cos(x^2).
    \]
    \begin{enumerate}
        \item Obtenir la dérivée première $f'(x)$ et identifier tous les points qui pourraient être un minimum ou un maximum local.
        \item Obtenir la dérivée seconde $f''(x)$.
        \item Réaliser un tableau reprenant les zones de croissance et de décroissance ainsi que les extrema de $f$ sur $[0, 2\sqrt{\pi}]$.
    \end{enumerate}

    \item \textbf{(UCLouvain -- Juillet 2023 -- Série 2, Q1 -- Étude complète)} On considère la fonction $f : \R \to \R$ définie par :
    \[
        f(x) = 3\cos(x) - \cos^3(x).
    \]
    \begin{enumerate}
        \item Obtenir $f'(x)$, identifier les candidats extrema locaux et trouver les racines de $f$.
        \item Obtenir $f''(x)$ et ses racines.
        \item Réaliser le tableau complet de variations et de concavité de $f$ sur $[0, 3\pi]$.
    \end{enumerate}

    \item \textbf{(UCLouvain -- Septembre 2023, Q1 -- Étude complète)} On considère la fonction $f : \D \subset \R \to \R$ définie par :
    \[
        f(x) = \frac{x}{\cos^2(x)}.
    \]
    Déterminer son domaine $\D$, sa dérivée première $f'(x)$, le tableau des valeurs de $f$ et $f'$ en $x \in \left\{0, \frac{\pi}{4}, \frac{\pi}{2}, \frac{3\pi}{4}, \pi\right\}$, et justifier pourquoi $f$ possède un minimum local sur $\left]\frac{3\pi}{4}, \pi\right[$.

    \item \textbf{(UCLouvain -- Juillet 2026, Q1 -- Étude complète)} On considère la fonction $f : \D \subset \R \to \R$ définie par :
    \[
        f(x) = (1 - x^2)\ln(1 - x) + x.
    \]
    \textbf{Aide :} $\lim_{y \to 0^+} y\ln(y) = 0$, $\ln(2) \simeq 0,6931$ et $e^{-1/2} \simeq 0,6065$.
    Déterminer $\D$, $\lim_{x \to 1} f(x)$, la dérivée première $f'(x)$ et ses deux racines $x_1, x_2$, la dérivée seconde $f''(x)$, puis dresser le tableau complet de variations et de concavité sur $[-1, 1[$.
\end{enumerate}
